\chapter{System Design}
\label{ch:system}

The decision models of \cref{ch:models} are only as trustworthy as the system that executes them
and the simulator that measures them. This chapter describes that system and the design choices
that make its numbers believable.

\section{Architecture}
\label{sec:architecture}

The system is a set of ROS~2 nodes around a Gazebo simulation (\cref{fig:architecture}): a
\textbf{generator} that turns one layout description into the simulated world, the navigation map,
the named poses and the maze distances; a \textbf{factory node} running the production model of the
three sections; a \textbf{mission executor} that turns one action into navigation goals, dwell
times and a report, and recovers when navigation fails; and a \textbf{decision service} that
answers \emph{what next?} and can be stopped or replaced without touching the robot.

\begin{figure}[H]
    \centering
    \begin{tikzpicture}[
        x=1.05cm, y=1.05cm, font=\small,
        box/.style={draw, rounded corners=2pt, align=center, inner sep=4pt,
                    minimum height=1.0cm, text width=2.5cm},
        src/.style={box, fill=black!5},
        gen/.style={box, fill=black!10},
        art/.style={box, fill=black!3, minimum height=0.6cm, font=\footnotesize,
                    text width=2.5cm},
        ros/.style={box, fill=black!8},
        ext/.style={box, dashed, fill=white},
        link/.style={-{Stealth[length=2mm]}, thick},
        thin/.style={-{Stealth[length=1.7mm]}, draw=black!60},
        obs/.style={-{Stealth[length=1.7mm]}, draw=black!60, dotted, thick},
        lbl/.style={font=\scriptsize, text=black!65, inner sep=1.5pt}]

        % ------------------------------------------------ generated once, before a run
        \node[src] (layout)    at (0,7)   {\texttt{layout.yaml}\\\footnotesize the hall as a grid};
        \node[gen] (generator) at (3.7,7) {\textbf{generator}};
        \node[art] (poses)     at (7.6,7.8) {poses $+$ distances};
        \node[art] (world)     at (7.6,7.0) {simulated world};
        \node[art] (map)       at (7.6,6.2) {navigation map};
        \draw[link] (layout) -- (generator);
        \draw[thin] (generator.east) -- (poses.west);
        \draw[thin] (generator.east) -- (world.west);
        \draw[thin] (generator.east) -- (map.west);

        % ---------------------------------------------------------- the running system
        \node[ros] (factory)  at (0,3.6)   {\textbf{factory node}\\\footnotesize sections A, B, C};
        \node[ros] (executor) at (3.7,3.6) {\textbf{mission}\\\textbf{executor}};
        \node[ext] (service)  at (11.6,3.6) {\textbf{decision service}\\\footnotesize the policy of Ch.~\ref{ch:models}};
        \node[ros] (state)    at (0,1.4)   {\textbf{robot state}\\\footnotesize battery, heat, wear};
        \node[ros] (nav)      at (3.7,1.4) {\textbf{Nav2}\\\footnotesize plan, follow, recover};
        \node[ros] (gazebo)   at (7.6,1.4) {\textbf{Gazebo}\\\footnotesize physics, sensors};
        \node[ext] (dash)     at (12.1,1.4) {\textbf{dashboard}\\\footnotesize read-only};

        \draw[thin] ([yshift=3pt]factory.east) -- node[lbl, above] {status}
                    ([yshift=3pt]executor.west);
        \draw[thin] ([yshift=-5pt]executor.west) -- node[lbl, below] {pickup}
                    ([yshift=-5pt]factory.east);
        \draw[link] ([yshift=3pt]executor.east) -- node[lbl, above] {state}
                    ([yshift=3pt]service.west);
        \draw[link] ([yshift=-5pt]service.west) -- node[lbl, below] {action}
                    ([yshift=-5pt]executor.east);
        \draw[link] (executor) -- node[lbl, right] {goal} (nav);
        \draw[link] (nav) -- node[lbl, above] {drive} (gazebo);
        \draw[thin] (state.north) -- node[lbl, left] {telemetry} (factory.south);
        \draw[thin] (gazebo.south) to[out=250, in=290, looseness=0.55]
                    node[lbl, below, pos=0.5] {odometry, lidar} (state.south);

        \begin{scope}[on background layer]
            \node[draw, dashed, rounded corners, fit=(layout)(generator)(poses)(map),
                  inner sep=9pt, label={[lbl]above:generated once, before a run}] (gbox) {};
            \node[draw, dashed, rounded corners, fit=(factory)(state)(executor)(nav)(gazebo),
                  inner sep=11pt, label={[lbl]below:ROS 2 nodes and the simulator}] (rbox) {};
        \end{scope}
        \draw[obs] (rbox.east |- dash) -- node[lbl, above, pos=0.45] {observes} (dash.west);
        \draw[thin, dotted] (gbox.south) -- node[lbl, right]
              {world, map, poses and distances} (rbox.north);
    \end{tikzpicture}
    \caption{System architecture. Everything geometric derives from one file (top); the decision
             service is a separate process that exchanges plain data with the executor, and the
             dashboard only observes.}
    \label{fig:architecture}
\end{figure}

One principle runs through the design: \textbf{the decision layer never sees ROS}. A policy
receives plain data and returns an action, which is what allows the identical policy object to run
inside the fast simulator for thousands of training shifts and inside the live service on the
robot. If the two had separate implementations, every result in \cref{ch:results} would depend on
their agreeing, and that agreement would have to be proved rather than assumed.

\section{One description of the world}
\label{sec:single-source}

The hall is described once, as an ASCII grid in which each character is a half-metre cell: wall,
machine, pickup zone, delivery point, charging station or free floor. A generator turns that file
into the Gazebo world, the occupancy map used for navigation, the named poses and the matrix of
shortest maze paths between them. Maintaining a simulated world and a navigation map separately is
a known trap --- the two drift apart, and the robot plans in a hall that differs from the one it
drives in --- and deriving both from one description makes that error impossible. It also means the
distances in \cref{tab:distances}, used by the fast simulator and the symbolic rules, are a
property of the simulated world rather than an approximation of it.

The first version of the hall demonstrated why this matters. It contained a \SI{0.5}{\metre} gap
between two wall segments; Nav2 treats only about \SI{0.22}{\metre} around an obstacle as lethal, so
its planner routed through the gap while the controller could not fit the robot through it. The
robot planned \SI{10.47}{\metre} from B to C and then drove \SI{24.4}{\metre} after re-routing. The
layout was widened, but the lasting outcome is the check added to the generator: it compares the
shortest path through comfortable space with the shortest path through any space wider than
\SI{0.25}{\metre}, and refuses to generate a world if the second is more than \SI{5}{\percent}
shorter for any pair of points, which is exactly the signature of a gap the robot cannot use. Fed
the old layout, it flags precisely the three pairs Nav2 had got wrong.

\section{Models shared by both simulators}
\label{sec:models-shared}

The factory model of \cref{sec:factory} and the robot model of \cref{sec:robot} are pure Python
with seeded random generators, and are used unchanged by the fast simulator and by the ROS nodes,
so a scenario behaves identically in both. Two implementation decisions are worth recording.

The three sections run inside a \emph{single} ROS node rather than one node each. This began as a
performance fix --- three Python nodes on the simulation clock, each waking on every clock message,
starved the navigation stack during start-up --- and the same investigation set the physics step to
\SI{4}{\milli\second}. It is the reason a full shift runs reliably on a laptop, which is what made
the Gazebo validation shifts possible at all.

The robot model is checked rather than trusted. A test asserts that its parameters imply physically
plausible magnitudes: consumption between \SIrange{5}{40}{\watt\hour\per\kilo\metre}, continuous
driving between one and four hours, a full charge between thirty minutes and two hours. That test
caught the first calibration error in this work --- a drive energy of
\SI{0.012}{\watt\hour\per\metre}, which reads plausibly but implies \SI{16.8}{\watt} of motor power
for a \SI{4.3}{\kilo\gram} robot and a runtime under an hour. The corrected value,
\SI{0.005}{\watt\hour\per\metre}, is the one used throughout.

\section{Executing an action}
\label{sec:executor}

The mission executor is the only component that touches both worlds. For a pickup it sends the
section's pose as a navigation goal, waits for arrival, stands still for the loading time on the
simulation clock, asks the section to hand over as many units as still fit, and records what the
action cost --- duration, distance, energy, battery --- next to what the model predicted it would
cost. That pairing is what later allows the fast simulator to be validated action by action.

Over a full shift a robot in a maze will fail to reach a goal: a costmap holds a stale obstacle, the
robot wedges against a wall, a planner refuses. Failure is therefore treated as part of the job and
escalated in fixed steps --- retry, clear both costmaps and retry, back up \SI{0.3}{\metre} and
retry, and finally abandon the goal and return to the charging station. Only if the charger is
unreachable as well does the robot halt and end the mission as \emph{navigation stuck}, the one
outcome that requires a person. Each drive also carries a timeout of three times its predicted
duration, so a robot that oscillates cannot hang the shift; \cref{sec:planner-study} contains a
shift in which this ladder ran to its end.

Arrival is verified rather than believed. Nav2 can report success without the robot being at the
goal: when a goal pose is blocked, its planner plans to the nearest free position within its
tolerance and the controller duly arrives \emph{there}. Every reported success is therefore checked
against the robot's localised position, and a deviation beyond \SI{0.35}{\metre} counts as a
failure. Without this check a blocked pickup point would silently become a successful pickup of
nothing, and the resulting data would be wrong in a way no later analysis could detect.

\section{Decision service and monitoring}
\label{sec:service}

The decision model runs as its own process and answers over HTTP: the executor sends the state and
the legal actions, the service returns the chosen action with its scores and explanation. The
separation means the model can be replaced without restarting the robot, that ``the AI is
unavailable'' is a state which can be produced deliberately by stopping the process, and that the
latency reported in \cref{sec:latency} is measured on the path the robot actually uses. If the
service does not answer, the executor falls back to the rule-based policy of
\cref{sec:rule-policy} and continues the shift, recording for each decision whether it came from
the model or the fallback: the failure degrades the quality of decisions rather than the
availability of the robot.

A web dashboard shows the live state --- sections, robot, shift totals, the maze with the robot's
position, and every decision with its explanation --- and is deliberately read-only, with no route
that accepts a command. Its map is drawn from the same layout description as the simulated world,
so what an operator sees cannot drift from where the robot is.

\section{Two simulators}
\label{sec:two-simulators}

The system exists twice. The full simulation --- Gazebo physics, Nav2, localisation, the executor
--- is faithful but runs in real time: one shift takes an hour. The fast simulator replaces driving
with the path matrix and the robot model, advances the same factory model in one-second steps, and
completes a shift in milliseconds. Training the scorer and evaluating nine policies over \num{3240}
shifts is possible only in the second; claiming that those results describe the first requires
evidence. That evidence is produced by replaying a completed Gazebo mission action by action in the
fast simulator: for the reference mission the totals agree to \SI{0.6}{\percent} in energy, with the
worst single action differing by \SI{10}{\percent} in duration. \Cref{sec:gazebo-validation} extends
this from one scripted mission to complete autonomous shifts under both decision models.
